\begin{abstract}
Multi-hop retrieval-augmented reasoning requires more than retrieving relevant passages: a system must decide which unresolved inference deserves the next unit of computation, when evidence should revise earlier beliefs, and when further work is no longer justified. We introduce \method, a training-free framework that represents test-time reasoning as a dynamic, typed hypergraph and uses that state to allocate language-model calls, tokens, retrievals, and verification effort. Evidence, claims, subgoals, and answer candidates retain separate belief channels rather than being prematurely compressed into one probability. Typed diffusion propagates uncertainty and downstream influence; explicit multi-premise joins create auditable intermediate claims; event-triggered editing repairs contradictions; and a multi-resource expected-value-of-computation controller selects both an operation and its fidelity. Safe termination distinguishes supported answers, abstention, and budget exhaustion. We evaluate \method{} on three multi-hop question-answering benchmarks under matched-compute protocols, with graph-memory, iterative-retrieval, and adaptive-RAG baselines. \tbd{Insert final aggregate quality, Pareto, revision, and reliability results after the frozen evaluation.} Development diagnostics indicate that the framework can produce non-uniform allocation and auditable three- and four-hop joins, while also exposing the remaining challenge: additional retrieval does not reliably translate into answer-relevant evidence. \method{} reframes test-time scaling as a state-dependent resource-allocation problem over explicit reasoning structure.
\end{abstract}
